\documentclass[12pt,openright,oneside,a4paper,english,brazil]{abntex2}
\usepackage[T1]{fontenc}
\usepackage[utf8]{inputenc}
\usepackage[alf]{abntex2cite}
\usepackage{lmodern,indentfirst,microtype,amsmath,amsfonts,amssymb,graphicx,xurl,xcolor,hyperref,booktabs}
\usepackage{tikz}
\usetikzlibrary{positioning,arrows.meta}

\titulo{Glossário de Conceitos Fundamentais de CV-QKD}
\autor{Willian Vinícius Oliveira Araújo}
\local{Campina Grande, PB, Brasil}
\data{Outubro de 2026}
\instituicao{
 Universidade Federal de Campina Grande (UFCG)
 \par
 Unidade Acadêmica de Engenharia Elétrica (UAEE)}
\tipotrabalho{Documentação Técnica -- PIBITI}
\preambulo{Glossário introdutório dos fundamentos utilizados no simulador de reconciliação de informação em sistemas CV-QKD.}

\definecolor{blue}{RGB}{41,5,195}
\makeatletter
\hypersetup{
 pdftitle={\@title},
 pdfauthor={\@author},
 pdfsubject={PIBITI - Reconciliação da informação em CV-QKD},
 pdfcreator={LaTeX with abnTeX2},
 pdfkeywords={PIBITI, CV-QKD, BI-AWGN, reconciliação},
 colorlinks=true,
 linkcolor=blue,
 citecolor=blue,
 filecolor=magenta,
 urlcolor=blue,
 bookmarksdepth=4
}
\makeatother
\setlength{\parindent}{1.3cm}
\setlength{\parskip}{0.2cm}
\begin{document}
\selectlanguage{brazil}
\frenchspacing
\imprimircapa
\textual

\chapter{Objetivo}
Reunir os principais conceitos utilizados no estudo inicial de um simulador elementar do canal BI-AWGN, conforme o plano de trabalho PIBITI 2026--2027. Este documento é introdutório e não descreve toda a cadeia física ou a segurança de um protocolo CV-QKD.

\chapter{Conceitos fundamentais}
\section{QKD -- Distribuição Quântica de Chaves}
\textit{Quantum Key Distribution}. Conjunto de protocolos destinados ao estabelecimento de chaves criptográficas entre participantes, com etapas quânticas e clássicas.

\section{CV-QKD -- QKD por Variáveis Contínuas}
\textit{Continuous-Variable Quantum Key Distribution}. Modalidade de QKD que utiliza observáveis contínuos, como quadraturas do campo eletromagnético. Alice e Bob obtêm dados correlacionados, mas não necessariamente idênticos.

\section{Reconciliação de informação}
Etapa de pós-processamento clássico destinada a corrigir diferenças entre dados correlacionados de Alice e Bob. Pode utilizar códigos de correção de erros, como LDPC.

\section{Bit e símbolo antipodal}
Um bit assume os valores 0 ou 1. No mapeamento antipodal adotado neste simulador, usa-se
\begin{equation}
0\mapsto +1,\qquad 1\mapsto -1,\qquad x=1-2b.
\end{equation}
Os símbolos $+1$ e $-1$ são opostos e permitem representar uma entrada binária para o canal.

\section{Variável aleatória e distribuição gaussiana}
Uma variável aleatória representa o resultado numérico de um fenômeno aleatório. A distribuição normal é descrita por média $\mu$ e variância $\sigma^2$:
\begin{equation}
Z\sim\mathcal{N}(\mu,\sigma^2).
\end{equation}
No simulador básico, adota-se $\mu=0$.

\section{Média, variância e desvio-padrão}
A média indica o centro da distribuição; a variância, a dispersão quadrática em torno da média; e o desvio-padrão é $\sigma=\sqrt{\sigma^2}$. Média zero não significa ausência de ruído.

\section{Canal BI-AWGN}
\textit{Binary-Input Additive White Gaussian Noise}. Modelo com entrada binária antipodal e ruído gaussiano branco aditivo:
\begin{equation}
Y=X+Z,\qquad X\in\{-1,+1\},\qquad Z\sim\mathcal{N}(0,\sigma^2).
\end{equation}
No modelo discreto empregado, as amostras de ruído são independentes e identicamente distribuídas.

\section{Detecção por decisão dura}
Para símbolos igualmente prováveis e ruído gaussiano simétrico, a estimativa de bit é obtida pelo sinal do valor recebido:
\begin{equation}
\hat b=\begin{cases}
0,&Y\ge 0,\\
1,&Y<0.
\end{cases}
\end{equation}

\section{BER -- Taxa de erro de bits}
\textit{Bit Error Rate}. Proporção de bits detectados incorretamente em uma simulação:
\begin{equation}
\mathrm{BER}=\frac{N_{\mathrm{erros}}}{N_{\mathrm{bits}}}.
\end{equation}
Para a entrada antipodal unitária do modelo, a probabilidade teórica de erro é $P_b=Q(1/\sigma)$, sob decisão ótima com probabilidades a priori iguais.

\section{LDPC -- Código de paridade de baixa densidade}
\textit{Low-Density Parity-Check}. Família de códigos corretores de erros descritos por matrizes de verificação de paridade esparsas. Sua implementação será tratada em etapas posteriores do projeto.

\chapter{Referências orientadoras}
\begin{itemize}
\item Plano de Trabalho PIBITI 2026--2027, INCT NanoQuantSS, seções 7--10.
\item MACKAY, D. J. C. \textit{Information Theory, Inference, and Learning Algorithms}. Cambridge University Press, 2003.
\item LEVERRIER, A. et al. \textit{Multidimensional reconciliation for a continuous-variable quantum key distribution}. Physical Review A, 77, 042325, 2008.
\item WEEDBROOK, C. et al. \textit{Gaussian quantum information}. Reviews of Modern Physics, 84, 621--669, 2012.
\end{itemize}

\end{document}
